% ==================================================================
% id:       cond-015
% topic:    הסתברות מותנית ואי-תלות
% subtopic: נוסחת ההסתברות השלמה ונוסחת בייס
% solution: full
% ==================================================================
אחרי האסון בכור האטומי בצ׳רנוביל הוחלט לבדוק קרינה רדיואקטיבית באזורים שונים. בכל אזור, ההסתברות לקרינה עודפת היא 0.02. המכשיר מראה תוצאה חיובית בהסתברות 0.95 אם יש קרינה עודפת, ובהסתברות 0.01 אם אין.
\begin{enumerate}
\item מה ההסתברות שהמכשיר יראה תוצאה חיובית?
\begin{solbox}
נסמן $R$ = „יש קרינה עודפת”, $+$ = „תוצאה חיובית”. לפי נוסחת ההסתברות השלמה
\[
\PP(+)=0.95\cdot0.02+0.01\cdot0.98=0.019+0.0098=\ans{0.0288}.
\]
\end{solbox}
\item באזור מסוים המכשיר הראה תוצאה שלילית. מה ההסתברות שבכל זאת יש שם קרינה עודפת?
\begin{solbox}
$\PP(-)=1-0.0288=0.9712$, ו-$\PP(-\mid R)=0.05$. לפי בייס
\[
\PP(R\mid -)=\frac{0.05\cdot0.02}{0.9712}=\frac{0.001}{0.9712}\approx\ans{0.00103}.
\]
מקרה כזה, שבו המכשיר מראה תוצאה שלילית למרות שיש קרינה עודפת, נקרא \textbf{תוצאה שלילית שגויה} (\textenglish{false negative}). המקרה ההפוך, תוצאה חיובית כשאין קרינה עודפת, נקרא \textbf{תוצאה חיובית שגויה} (\textenglish{false positive}); בשאלה הזו ההסתברות שלה היא 0.01.
\end{solbox}
\end{enumerate}
